\kbvsechead{(2)\hspace{0.45em}ಸರ್ವನಾಮಗಳು :}

\indent “ದೇವರು ದೊಡ್ಡವನು ; ದೇವರು ಸರ್ವಶಕ್ತ” — ಈ ವಾಕ್ಯಗಳನ್ನು ಓದಿ ನೋಡಿ. ಎರಡನೆಯ ವಾಕ್ಯದಲ್ಲಿ “ದೇವರು” ಎಂಬ ನಾಮಪದ ಪುನರಾವೃತ್ತವಾಗಿದೆ. ಈ ಪದಕ್ಕೆ ಬದಲಾಗಿ ಅವನು ಎಂಬ ಪದ ಬಳಸಿ ದಲ್ಲಿ ಓದಲು ಹಿತವೆನಿಸುವುದು. “ದೇವರು ದೊಡ್ಡವನು; ಅವನು ಸರ್ವಶಕ್ತ.”\par

\noindent\hspace*{2.2em}ಇದೇ ರೀತಿಯ ಹಲವು ಉದಾಹರಣೆಗಳನ್ನು ಪರಿಶೀಲಿಸೋಣ.\par

\vspace{2.5pt plus 3pt minus 0.5pt}\noindent\hspace*{1.2em}\kbvnum{(1)}\hspace{0.55em}ಪಾರ್ವತಿಯು ಶಿವನ ಹೆಂಡತಿ; ಅವಳು ಪರ್ವತೇಶ್ವರನ ಮಗಳು.\par

\vspace{2.5pt plus 3pt minus 0.5pt}\noindent\hspace*{1.2em}\kbvnum{(2)}\hspace{0.55em}ಬೆಂಗಳೂರು ದೊಡ್ಡ ಪಟ್ಟಣ; ಅದು ಕರ್ನಾಟಕದ ರಾಜಧಾನಿ.\par

\vspace{2.5pt plus 3pt minus 0.5pt}\noindent\hspace*{1.2em}\kbvnum{(3)}\hspace{0.55em}ನಾನು ಇಲ್ಲಿಯೇ ಕಾದಿರುತ್ತೇನೆ ; ನೀನು ಆದಷ್ಟು ಬೇಗ ಬಾ.\par

\vspace{2.5pt plus 3pt minus 0.5pt}\noindent\hspace*{1.2em}\kbvnum{(4)}\hspace{0.55em}ಇವರು ಯಾರು ?...{[}ಈ ಪ್ರಶ್ನೆಗೆ ಉತ್ತರವಾಗಿ “ಇವರು ಗಾಂಧೀಜಿ” ಎಂದು ಹೇಳುವುದಾದರೆ, ಗಾಂಧೀಜಿ. ಎಂಬ ನಾಮಪದಕ್ಕೆ ಬದ ಲಾಗಿ ಯಾರು ಎಂಬ ಪ್ರಶ್ನಾರ್ಥಕ ಪದ ಬಳಸಿರುವುದು ವ್ಯಕ್ತವಾಗುವುದು.{]}\par

\indent ಮೇಲಿನ ಉದಾಹರಣೆಗಳಲ್ಲಿ ನಾಮಪದಗಳಿಗೆ ಬದಲಾಗಿ ಅವನು ಅವಳು, ಅದು, ನಾನು, ನೀನು, ಯಾರು ಎಂಬ ಪದಗಳನ್ನು ಬಳಸಿರುವುದು ಕಂಡುಬರುತ್ತದೆ. ಹೀಗೆ ನಾಮಪದಗಳಿಗೆ ಪ್ರತಿಯಾಗಿ ಅವುಗಳ ಸ್ಥಾನದಲ್ಲಿ ಬಳಸುವ ಪದಗಳಿಗೆ “ಸರ್ವನಾಮ” ಗಳೆಂದು ಹೆಸರು.\par

\kbvsechead{ಸರ್ವನಾಮಗಳಲ್ಲಿ ಮೂರು ವಿಧ :}

\vspace{2.5pt plus 3pt minus 0.5pt}\noindent\hspace*{1.2em}\kbvnum{(1)}\hspace{0.55em}ಪುರುಷಾರ್ಥಕ ಸರ್ವನಾಮಗಳು\par

\vspace{2.5pt plus 3pt minus 0.5pt}\noindent\hspace*{1.2em}\kbvnum{(2)}\hspace{0.55em}ಆತ್ಮಾರ್ಥಕ ಸರ್ವನಾಮಗಳು\par

\vspace{2.5pt plus 3pt minus 0.5pt}\noindent\hspace*{1.2em}\kbvnum{(3)}\hspace{0.55em}ಪ್ರಶ್ನಾರ್ಥಕ ಸರ್ವನಾಮಗಳು\par

\kbvsechead{(1)\hspace{0.45em}ಪುರುಷಾರ್ಥಕ ಸರ್ವನಾಮಗಳು :}

\indent ಉದಾಹರಣೆಗೆ ನಾನು, ನೀನು, ಅವನು, ‘ಇವನು’, ‘ಅದು’, ‘ಇದು’ ಇತ್ಯಾದಿ ಪುರುಷಾರ್ಥಕ ಸರ್ವನಾಮಗಳು. ಈ ಸರ್ವನಾಮಗಳನ್ನು ಮೂರು ಗುಂಪುಗಳಾಗಿ ವಿಂಗಡಿಸಬಹುದು.\par
